\chapter{DSL Implementation}

The implementation of the DaQuVa DSL follows a classic compiler pipeline: lexer $\rightarrow$ parser $\rightarrow$ AST. All components are written in pure Python~3 and are directly derived from the formal BNF grammar presented in Chapter~4. The code is kept minimal, readable, and fully aligned with the safety-first and readability principles of the language.

\section{Lexer}

The lexer is a hand-written, regex-based tokenizer that converts the raw source code into a flat stream of tokens. It explicitly recognises every terminal symbol defined in the BNF (keywords, identifiers, literals, operators, punctuation) while discarding insignificant whitespace and comments. Although the grammar now defines <\texttt{\<SingleSpace\>}> and <\texttt{\<OneOrMoreSpaces\>}> for formal completeness, the lexer follows the standard convention (explicitly stated in the BNF note) and treats whitespace purely as a token separator.

Key design decisions:
\begin{itemize}
    \item All reserved words are turned into a single \texttt{KW} token.
    \item Whitespace (\texttt{[ \t\r\n]+}) and single-line comments (\texttt{\#.*}) are silently skipped.
    \item Multi-character operators (\texttt{->}, \texttt{==}, \texttt{!=}, \texttt{>=}, \texttt{<=}) are recognised before single-character ones.
    \item String literals have their surrounding quotes stripped.
    \item Numbers are immediately converted to Python \texttt{int}.
\end{itemize}

\begin{lstlisting}[style=dsl, caption={Core lexer implementation (excerpt)}, label={lst:lexer}]
KEYWORDS = frozenset({
    'my', 'db', 'function', 'return', 'end', 'scan', 'filter', 'merge',
    'output', 'newTable', 'findLocalDuplicates', 'where', 'AND', 'OR',
    'starts_with', 'contains', 'ends_with', 'sort', 'by', 'asc', 'desc',
    'limit', 'summarize',
    'columns', 'using', 'tool', 'console', 'file', 'from', 'addColumns',
    'Here', 'allowDanger'
})

class Lexer:
    _RULES = [
        ('COMMENT',   r'#[^\n]*'),
        ('STRING',    r'"[^"]*"'),
        ('NUMBER',    r'\d+'),
        ('ARROW',     r'->'),
        ('COMP_OP',   r'!=|==|>=|<=|>|<'),
        ('COLON',     r':'),
        ('ASSIGN',    r'='),
        ('SEMICOLON', r';'),
        ('LPAREN',    r'\('),
        ('RPAREN',    r'\)'),
        ('LBRACKET',  r'\['),
        ('RBRACKET',  r'\]'),
        ('DOT',       r'\.'),
        ('COMMA',     r','),
        ('WORD',      r'[A-Za-z_]\w*'),
        ('SKIP',      r'[ \t\r\n]+'),
        ('MISMATCH',  r'.'),
    ]

    def __init__(self, source: str):
        self.tokens = []
        self.pos = 0
        for mo in re.finditer('|'.join(f'(?P<{n}>{p})' for n,p in self._RULES), source):
            kind = mo.lastgroup
            if kind in ('SKIP', 'COMMENT'):
                continue
            if kind == 'MISMATCH':
                raise LexerError(...)
            if kind == 'STRING':
                self.tokens.append(Token('STRING', mo.group()[1:-1], mo.start()))
            elif kind == 'NUMBER':
                self.tokens.append(Token('NUMBER', int(mo.group()), mo.start()))
            elif kind == 'WORD':
                k = 'KW' if mo.group() in KEYWORDS else 'ID'
                self.tokens.append(Token(k, mo.group(), mo.start()))
            else:
                self.tokens.append(Token(kind, mo.group(), mo.start()))

    def peek(self, offset=0):
        idx = self.pos + offset
        return self.tokens[idx] if idx < len(self.tokens) else Token('EOF', None, -1)

    def consume(self, kind=None, value=None):
        tok = self.peek()
        if kind and tok.kind != kind or value and tok.value != value:
            raise SyntaxError(...)
        self.pos += 1
        return tok
\end{lstlisting}

The lexer produces exactly the token stream expected by the recursive-descent parser and handles all edge cases from the illustrative example in Section~4.4.

\section{Parser}

The parser is a classic recursive-descent implementation that mirrors the BNF structure one-to-one. It uses the lexer’s \texttt{peek}/\texttt{consume} interface and builds the AST on the fly. Two small practical adaptations were introduced for real-world usability while staying faithful to the formal grammar:

\begin{itemize}
    \item Function bodies are terminated by the keyword \texttt{end} (the original BNF used pure recursion without an explicit closer; \texttt{end} makes the language non-indentation-sensitive).
    \item The \texttt{return} statement (already present in the keyword table) was given its own AST node.
\end{itemize}

\begin{lstlisting}[style=dsl, caption={Parser entry point and statement dispatcher (excerpt)}, label={lst:parser}]
class Parser:
    def parse(self) -> Program:
        stmts = []
        while not self._at_end():
            stmts.append(self._parse_statement())
        return Program(stmts)

    def _parse_statement(self):
        tok = self._peek()
        if tok.kind == 'KW':
            if tok.value == 'my':          return self._parse_db_init()
            if tok.value == 'function':    return self._parse_function_def()
            if tok.value == 'return':      return self._parse_return()
            node = self._parse_command()
            self._consume('SEMICOLON')
            return node
        if tok.kind == 'ID':
            if self._is('ASSIGN', offset=1):
                return self._parse_assignment()
            if self._is('LPAREN', offset=1):
                node = self._parse_function_call()
                self._consume('SEMICOLON')
                return node
        raise ParseError(...)

    # Individual command parsers (scan, filter, findLocalDuplicates, merge,
    # output, newTable) follow the BNF productions exactly and are omitted
    # here for brevity. Each returns the corresponding AST node.
\end{lstlisting}

Every parsing method directly corresponds to a BNF production (e.g.\ \texttt{\_parse\_scan} implements \bnf{<scan\_command>}). Compound filters are parsed through separate \texttt{\_or\_condition} and \texttt{\_and\_condition} methods, which gives \texttt{AND} higher precedence than \texttt{OR}. The parser also recognises the latest command forms for \texttt{sort}, \texttt{limit}, \texttt{contains}, and \texttt{ends\_with}, and accepts a direct function call after \texttt{return}. Error messages include source position for excellent developer experience.

\section{AST}

The Abstract Syntax Tree is defined with Python \texttt{dataclass}es, giving a clean, immutable, and type-safe representation that exactly mirrors the BNF non-terminals. All nodes inherit from a common \texttt{ASTNode} base class.

\begin{lstlisting}[style=dsl, caption={Complete AST node hierarchy}, label={lst:ast}]
@dataclass
class ASTNode: pass

@dataclass
class Program(ASTNode):
    statements: List[ASTNode]

@dataclass
class DatabaseInit(ASTNode):
    connection: str

@dataclass
class FunctionDef(ASTNode):
    name: str
    params: List[str]
    body: List[ASTNode]

@dataclass
class Assignment(ASTNode):
    name: str
    value: ASTNode

@dataclass
class ReturnStatement(ASTNode):
    value: ASTNode

@dataclass
class Scan(ASTNode):
    table: ASTNode          # Literal | Variable
    column: str
    tool: str

@dataclass
class Filter(ASTNode):
    source: str
    condition: ASTNode

@dataclass
class CompoundCondition(ASTNode):
    left: ASTNode
    op: str
    right: ASTNode

@dataclass
class FindDuplicates(ASTNode):
    source: str
    columns: List[str]
    tool: Optional[str]

@dataclass
class Merge(ASTNode):
    source: str
    threshold: Optional[int]

@dataclass
class Output(ASTNode):
    source: str
    target: str
    path: Optional[str] = None
    threshold: Optional[int] = None

@dataclass
class NewTable(ASTNode):
    name: str
    source: Optional[str]
    columns: List[ASTNode]

@dataclass
class SortExpression(ASTNode):
    source: str
    column: str
    ascending: bool

@dataclass
class LimitExpression(ASTNode):
    source: str
    count: int

@dataclass
class SummaryExpression(ASTNode):
    source_names: Tuple[str, ...]   # one or more stages to summarize

@dataclass
class BinaryOp(ASTNode):
    left: ASTNode
    op: str
    right: ASTNode

@dataclass
class FunctionCall(ASTNode):
    name: str
    args: List[ASTNode]

@dataclass
class Variable(ASTNode):
    name: str

@dataclass
class Literal(ASTNode):
    value: Union[int, str]
\end{lstlisting}

A simple \texttt{print\_ast} pretty-printer (included in the implementation) produces human-readable tree views for debugging. The AST serves as the single source of truth for all subsequent phases (semantic analysis, execution engine, pretty-printing, etc.).

The lexer, parser, and AST together form a complete front-end that can already parse every valid DaQuVa program defined in the language specification.

\section{Semantic Analysis and Runtime}

The runtime phase resolves parsed AST nodes against the active execution
context. It registers CSV and SQLite connections, stores in-memory variables,
tracks user-defined functions, and turns table expressions into lazy logical
plans. This design keeps parsing independent from file systems and databases:
the parser only validates syntax, while the runtime validates names, columns,
tool availability, and output destinations.

Important runtime checks added in the current implementation include:

\begin{itemize}
    \item \textbf{Column validation}: filter conditions and projected columns
          are checked against the available schema. A misspelled condition
          column now raises an explicit error that lists available columns
          instead of silently producing an empty result.
    \item \textbf{Function return handling}: function bodies may return a
          variable, literal, or another function call, enabling composition of
          small reusable routines.
    \item \textbf{Controlled persistence}: output to a database destination
          accepts an \texttt{allowDanger} flag. The SQLite writer uses this
          flag when replacing persisted data, keeping potentially destructive
          writes visible in DSL scripts.
    \item \textbf{Tool aliases}: \texttt{levenshtein\_matcher} is resolved as
          an alias of the fuzzy duplicate detection implementation.
\end{itemize}

\section{Execution Engine}

The execution engine evaluates lazy logical plans and chooses the most suitable
backend strategy. For CSV-backed data, plans are materialized and evaluated in
Python. For SQLite-backed data, simple relational operations are compiled into
SQL whenever possible.

The latest implementation extends the executable plan set with
\texttt{SortPlan} and \texttt{LimitPlan}. In Python execution, \texttt{sort}
uses a mixed-type safe key so that empty values, numbers, booleans, and strings
do not produce unhandled \texttt{TypeError} exceptions. For SQLite plans,
\texttt{filter}, \texttt{sort}, and \texttt{limit} are pushed down into SQL:
\texttt{WHERE}, \texttt{ORDER BY}, and \texttt{LIMIT} are generated directly
against the source table.

String predicates are also executable in both backends. In memory, they are
evaluated with Python string operations. In SQLite, \texttt{starts\_with},
\texttt{contains}, and \texttt{ends\_with} are translated into parameterized
\texttt{LIKE} expressions:

\begin{lstlisting}[style=dsl]
email starts_with "a"    -> email LIKE "a%"
email contains "@gmail"  -> email LIKE "%@gmail%"
email ends_with ".org"   -> email LIKE "%.org"
\end{lstlisting}

Several correctness fixes were made in the execution layer. CSV input values
are stripped of surrounding whitespace before storage, preventing equality
filters from failing on visually identical strings. Numeric comparisons now
produce a clear \texttt{ValueError} when ordered comparisons are attempted
against incompatible non-numeric strings. SQLite table replacement was also
made safer by creating a temporary replacement table and atomically renaming it
only after creation succeeds.

The plan set is further extended with a \texttt{SummaryPlan}, which carries a
tuple of \texttt{SummarySource} entries (the stage name, its logical plan, and
its column list). When the engine encounters a \texttt{SummaryPlan} it
materialises each source and reduces it to one metrics row via a small
\texttt{\_summary\_row} helper. That helper counts distinct
\texttt{duplicate\_group\_id} values, sums the \texttt{merged\_from\_count}
annotation to recover how many input rows a merged result represents, and uses a
tolerant truthiness check for the \texttt{typo\_suspect} flag so that string,
boolean, and numeric encodings are all interpreted consistently. Because the
summary is emitted as an ordinary list of row dictionaries with a fixed
\texttt{SUMMARY\_COLUMNS} schema, it flows through the existing CSV and SQLite
writers unchanged.

\section{Interactive REPL}

DaQuVa now includes an interactive REPL started with:

\begin{lstlisting}[style=dsl]
python -m daquva --repl
\end{lstlisting}

The REPL keeps a persistent runtime context, so users can declare connections,
assign intermediate tables, inspect variables, and print results without
rerunning an entire script. It buffers multi-line input until a semicolon or
function-closing brace is present, then parses and executes the accumulated
snippet. Meta-commands provide lightweight inspection:

\begin{itemize}
    \item \texttt{:help} prints available REPL commands and examples.
    \item \texttt{:tables} lists active connections.
    \item \texttt{:vars} lists in-memory variables and table schemas.
    \item \texttt{:quit} exits the session.
\end{itemize}

After an assignment, the REPL reports the variable name, row count, column
count, and column names. This immediate feedback is useful for a demonstration
setting because each transformation can be inspected before the next command is
entered.

\section{Showcase Demo}

A new end-to-end showcase script was added under
\texttt{code/demo/demo\_showcase}. It demonstrates the latest implemented
features in one narrative pipeline: typo detection, compound
\texttt{AND}/\texttt{OR} filters, provider filtering with \texttt{contains},
top-$n$ exploration with \texttt{sort} and \texttt{limit}, domain filtering
with \texttt{ends\_with}, fuzzy duplicate detection and merge, CSV export, and
SQLite persistence with \texttt{allowDanger}.

To make the pipeline auditable, the demo input was enlarged to 30 customer rows
containing deliberate near-duplicates and email typos, and the script ends with
a \texttt{summarize} step that proves the cleaning effect:

\begin{lstlisting}[style=dsl]
proof = summarize raw, paid_adults, dupes, merged;
proof -> file "output/proof_summary.csv";
save proof into proof_summary allowDanger;
\end{lstlisting}

This closes the loop required of a data-quality tool: a concrete CSV
\emph{input}, cleaned CSV and SQLite \emph{outputs}, and an aggregate summary
that is itself written to both CSV (\texttt{output/proof\_summary.csv}) and a
SQLite table (\texttt{proof\_summary}). The summary does not merely append
columns; it reports before/after row counts so a reviewer can confirm, for
example, that 30 raw rows were reduced to 15 merged rows representing 25 inputs.

